% =============================================================================
% Rapport technique – Système d'Authentification par Reconnaissance Faciale
% Auteur : Fares Jebali
% Date   : Avril 2026
% =============================================================================
\documentclass[12pt,a4paper]{article}

% --- Encodage & langue -------------------------------------------------------
\usepackage[T1]{fontenc}
\usepackage[utf8]{inputenc}
\usepackage[french]{babel}

% --- Mise en page ------------------------------------------------------------
\usepackage[top=2.5cm,bottom=2.5cm,left=3cm,right=2.5cm]{geometry}
\usepackage{setspace}
\onehalfspacing

% --- Typographie -------------------------------------------------------------
\usepackage{lmodern}
\usepackage{microtype}

% --- Graphiques et couleurs --------------------------------------------------
\usepackage{xcolor}
\usepackage{graphicx}
\usepackage{tikz}
\usetikzlibrary{shapes.geometric, arrows.meta, positioning, fit, backgrounds}

% --- Tableaux ----------------------------------------------------------------
\usepackage{booktabs}
\usepackage{longtable}
\usepackage{array}
\usepackage{multirow}

% --- Maths -------------------------------------------------------------------
\usepackage{amsmath}
\usepackage{amssymb}

% --- Code source -------------------------------------------------------------
\usepackage{listings}
\usepackage{inconsolata}

% --- Liens -------------------------------------------------------------------
\usepackage[colorlinks=true,linkcolor=NavyBlue,urlcolor=NavyBlue,citecolor=NavyBlue]{hyperref}

% --- En-têtes ----------------------------------------------------------------
\usepackage{fancyhdr}
\pagestyle{fancy}
\fancyhf{}
\renewcommand{\headrulewidth}{0.4pt}
\fancyhead[L]{\small Système d'Authentification par Reconnaissance Faciale}
\fancyhead[R]{\small Fares Jebali}
\fancyfoot[C]{\thepage}

% --- Sections stylisées ------------------------------------------------------
\usepackage{titlesec}
\definecolor{TitleBlue}{RGB}{0,64,128}
\definecolor{AccentBlue}{RGB}{30,100,180}
\definecolor{LightGray}{RGB}{245,245,245}
\definecolor{NavyBlue}{RGB}{0,50,120}
\definecolor{CodeBg}{RGB}{248,248,248}

\titleformat{\section}
  {\large\bfseries\color{TitleBlue}}
  {\thesection}{1em}{}[\vspace{-0.2em}\color{TitleBlue}\hrule\vspace{0.3em}]

\titleformat{\subsection}
  {\normalsize\bfseries\color{AccentBlue}}
  {\thesubsection}{1em}{}

\titleformat{\subsubsection}
  {\normalsize\itshape\color{AccentBlue}}
  {\thesubsubsection}{1em}{}

% --- Style de code -----------------------------------------------------------
\lstdefinestyle{pycode}{
  language=Python,
  backgroundcolor=\color{CodeBg},
  basicstyle=\ttfamily\small,
  keywordstyle=\color{TitleBlue}\bfseries,
  stringstyle=\color{orange!70!black},
  commentstyle=\color{gray}\itshape,
  numbers=left,
  numberstyle=\tiny\color{gray},
  frame=single,
  framerule=0.4pt,
  rulecolor=\color{gray!40},
  breaklines=true,
  showstringspaces=false,
  tabsize=4,
}
\lstset{style=pycode}

% --- Boîtes colorées ---------------------------------------------------------
\usepackage{tcolorbox}
\tcbuselibrary{skins}
\newtcolorbox{infobox}[1]{
  colback=AccentBlue!8,
  colframe=AccentBlue!60,
  fonttitle=\bfseries,
  title=#1,
  arc=4pt,
}
\newtcolorbox{formulebox}[1]{
  colback=TitleBlue!6,
  colframe=TitleBlue!50,
  fonttitle=\bfseries,
  title=#1,
  arc=4pt,
}

% =============================================================================
\begin{document}
% =============================================================================

% ---- Page de titre ----------------------------------------------------------
\begin{titlepage}
  \centering
  \vspace*{1cm}

  {\Huge\bfseries\color{TitleBlue}
    Système d'Authentification\\[0.4em]
    par Reconnaissance Faciale\par}

  \vspace{0.6cm}
  {\large\color{gray} Rapport Technique Complet\par}

  \vspace{1.5cm}
  \begin{tikzpicture}
    \draw[TitleBlue,line width=2pt] (0,0) -- (12,0);
    \draw[AccentBlue,line width=1pt] (0,-0.3) -- (12,-0.3);
  \end{tikzpicture}

  \vspace{1.5cm}

  % Architecture overview box
  \begin{tcolorbox}[
    colback=LightGray,
    colframe=TitleBlue,
    width=0.85\textwidth,
    arc=6pt,
    title={\bfseries\color{white} Vue d'ensemble du système},
    coltitle=white,
    attach boxed title to top center={yshift=-3mm},
    boxed title style={colback=TitleBlue,arc=3pt},
  ]
    \centering
    \small
    \begin{tabular}{ll}
      \textbf{Langage} & Python 3.12 \\
      \textbf{Interface} & Tkinter (application desktop) \\
      \textbf{Deep Learning} & FaceNet, ArcFace (via DeepFace / TensorFlow) \\
      \textbf{Vision classique} & OpenCV, Dlib, scikit-learn \\
      \textbf{Sécurité} & AES-256-GCM, PBKDF2-SHA256, BioHashing \\
      \textbf{Détection anti-spoof} & EAR (blink), LBP texture, FFT deepfake \\
    \end{tabular}
  \end{tcolorbox}

  \vspace{2cm}

  {\large\bfseries Fares Jebali\par}
  \vspace{0.3cm}
  {\normalsize Avril 2026\par}

  \vfill
\end{titlepage}

% ---- Table des matières -----------------------------------------------------
\tableofcontents
\newpage

% =============================================================================
\section{Introduction}
% =============================================================================

La biométrie faciale est aujourd'hui l'une des modalités biométriques les plus
utilisées dans les systèmes d'authentification modernes, en raison de sa nature
non intrusive, de sa facilité de déploiement et de la richesse des informations
qu'elle fournit. Des applications telles que le déverrouillage de smartphones,
le contrôle d'accès physique et l'authentification bancaire ont popularisé cette
technologie.

Le présent projet implémente un \textbf{système d'authentification complet basé
sur la reconnaissance faciale}, conçu comme une application desktop. Il combine
des approches classiques (Eigenfaces, LBPH) et profondes (FaceNet, ArcFace),
une détection de vivacité pour résister aux attaques par présentation, une
détection de \textit{deepfakes}, et des mécanismes de protection de la vie
privée (biométrie annulable, chiffrement AES-256-GCM).

\begin{infobox}{Objectifs du système}
  \begin{itemize}
    \item Authentifier un utilisateur en temps réel via sa webcam.
    \item Résister aux attaques par présentation (photos, vidéos, deepfakes).
    \item Protéger les données biométriques par chiffrement et templates
          révocables.
    \item Fournir une interface graphique intuitive et sécurisée.
  \end{itemize}
\end{infobox}

% =============================================================================
\section{Architecture Générale}
% =============================================================================

\subsection{Structure des modules}

Le système est organisé en couches fonctionnelles indépendantes communicant via
des interfaces bien définies. La figure~\ref{fig:arch} illustre cette
organisation.

\begin{figure}[h!]
\centering
\begin{tikzpicture}[
  node distance=0.6cm and 1.2cm,
  box/.style={rectangle, rounded corners=4pt, minimum width=3.2cm,
              minimum height=0.8cm, align=center, font=\small\bfseries,
              draw, thick},
  layer/.style={rectangle, rounded corners=6pt, draw=gray!50, dashed,
                inner sep=8pt},
  arrow/.style={-Stealth, thick, AccentBlue},
]

% Couche UI
\node[box, fill=AccentBlue!20, draw=AccentBlue] (ui) {Interface\\Tkinter (app.py)};

% Couche Orchestration
\node[box, fill=TitleBlue!15, draw=TitleBlue, below=of ui] (enroll) {Enrollment\\Manager};
\node[box, fill=TitleBlue!15, draw=TitleBlue, right=1.2cm of enroll] (auth) {Authenticator\\(fusion)};
\node[box, fill=TitleBlue!15, draw=TitleBlue, left=1.2cm of enroll] (sec) {Security\\Monitor};

% Couche Traitement
\node[box, fill=orange!15, draw=orange!70, below=0.7cm of enroll] (pre) {Preprocessing\\(détection + alignement)};
\node[box, fill=orange!15, draw=orange!70, right=1.2cm of pre] (deep) {Deep Model\\FaceNet / ArcFace};
\node[box, fill=orange!15, draw=orange!70, left=1.2cm of pre] (class) {Classical Model\\Eigenfaces / LBPH};

% Couche Anti-spoof
\node[box, fill=red!12, draw=red!60, below=0.7cm of pre] (live) {Liveness\\(EAR + LBP)};
\node[box, fill=red!12, draw=red!60, right=1.2cm of live] (deep2) {Deepfake\\(FFT)};

% Couche Sécurité données
\node[box, fill=green!12, draw=green!60, below=0.7cm of live] (vault) {Vault\\AES-256-GCM};
\node[box, fill=green!12, draw=green!60, right=1.2cm of vault] (cancel) {Cancelable\\BioHashing};
\node[box, fill=green!12, draw=green!60, left=1.2cm of vault] (store) {User Store\\(JSON)};

% Flèches
\draw[arrow] (ui) -- (enroll);
\draw[arrow] (ui) -- (auth);
\draw[arrow] (ui) -- (sec);
\draw[arrow] (enroll) -- (pre);
\draw[arrow] (enroll) -- (deep);
\draw[arrow] (enroll) -- (class);
\draw[arrow] (auth) -- (deep);
\draw[arrow] (auth) -- (class);
\draw[arrow] (pre) -- (live);
\draw[arrow] (pre) -- (deep2);
\draw[arrow] (enroll) -- (vault);
\draw[arrow] (enroll) -- (cancel);
\draw[arrow] (ui) -- (store);

\end{tikzpicture}
\caption{Architecture modulaire du système d'authentification faciale.}
\label{fig:arch}
\end{figure}

\subsection{Description des modules}

\begin{longtable}{>{\bfseries}p{3.8cm} p{10cm}}
\toprule
Module & Rôle \\
\midrule
\endhead
\bottomrule
\endfoot
app.py & Interface graphique principale. Gère les écrans (accueil, inscription,
         connexion, bienvenue) et orchestre les appels aux autres modules. \\[4pt]
preprocessing.py & Détection de visage (Haar + HOG), extraction de 68 points
                   repères, alignement affine, normalisation et égalisation
                   d'histogramme. \\[4pt]
classical\_model.py & Modèles classiques : Eigenfaces (PCA + k-NN) et LBPH
                      (Local Binary Pattern Histograms). \\[4pt]
deep\_model.py & Extraction d'embeddings via FaceNet ou ArcFace (bibliothèque
                 DeepFace). Comparaison par similarité cosinus. \\[4pt]
authenticator.py & Fusion pondérée des scores classique et profond.
                   Décision binaire GRANTED / DENIED. \\[4pt]
enrollment.py & Inscription des utilisateurs : capture webcam ou images,
                extraction d'embeddings, stockage dans la galerie et le coffre. \\[4pt]
liveness.py & Détection de vivacité : comptage de clignements (EAR) et
              analyse de texture (LBP). \\[4pt]
deepfake\_detector.py & Détection de visages synthétiques par analyse spectrale
                        (FFT). \\[4pt]
cancelable\_biometrics.py & Protection des templates par projection orthogonale
                             aléatoire (BioHashing) liée à un token secret. \\[4pt]
vault.py & Coffre chiffré AES-256-GCM avec clé dérivée par PBKDF2-SHA256. \\[4pt]
security\_monitor.py & Limitation de débit (sliding window) et verrouillage
                       progressif après tentatives échouées. \\[4pt]
user\_store.py & Base de données JSON des profils utilisateurs
                 (nom, prénom, date de naissance). \\[4pt]
lighting.py & Normalisation de l'éclairage avant traitement. \\[4pt]
evaluation.py & Calcul des métriques FAR, FRR, EER et courbe ROC. \\[4pt]
config.py & Source unique de vérité pour tous les hyperparamètres et chemins. \\
\end{longtable}

% =============================================================================
\section{Technologies Utilisées}
% =============================================================================

\subsection{Langage et environnement}

Le projet est entièrement développé en \textbf{Python 3.12}, choisi pour la
richesse de son écosystème en vision par ordinateur et apprentissage automatique.

\subsection{Bibliothèques principales}

\begin{longtable}{>{\bfseries}p{4cm} p{3.5cm} p{6cm}}
\toprule
Bibliothèque & Version & Rôle \\
\midrule
\endhead
\bottomrule
\endfoot
OpenCV (contrib) & 4.9.0.80 & Détection Haar, LBPH, traitement d'image,
                               webcam, égalisation d'histogramme. \\[3pt]
Dlib & 19.24.4 & Détecteur HOG+SVM, prédicteur de 68 points repères faciaux,
                  alignement affine. \\[3pt]
DeepFace & 0.0.91 & Wrapper unifié pour FaceNet et ArcFace ;
                     extraction d'embeddings profonds. \\[3pt]
TensorFlow & 2.16.1 & Backend de calcul pour les modèles profonds
                       (FaceNet, ArcFace). \\[3pt]
scikit-learn & 1.4.2 & PCA (Eigenfaces), k-NN, StandardScaler,
                        LabelEncoder, Pipeline. \\[3pt]
NumPy & 1.26.4 & Calcul matriciel, manipulation des embeddings,
                  FFT, projections orthogonales. \\[3pt]
SciPy & 1.13.0 & Distance cosinus, distance euclidienne (EAR). \\[3pt]
cryptography & 42.0.5 & AES-256-GCM (AESGCM), PBKDF2HMAC, hachages
                         cryptographiques. \\[3pt]
Pillow & 10.3.0 & Chargement et conversion d'images pour l'interface. \\[3pt]
Matplotlib / Seaborn & 3.8.4 / 0.13.2 & Tracé des courbes ROC, FAR/FRR
                                          pour l'évaluation. \\[3pt]
Pandas & 2.2.2 & Export des métriques d'évaluation en CSV. \\[3pt]
Tkinter & stdlib & Interface graphique desktop (inclus dans Python). \\
\end{longtable}

\subsection{Organisation des données sur disque}

\begin{lstlisting}[language={},caption={Arborescence des données persistées}]
face_auth_system/
|-- data/
|   |-- users.json          # Profils utilisateurs (nom, dob)
|   |-- enrolled_faces/     # Recadrages pour renentrainement classique
|   |-- logs/               # Journaux d'authentification
|   |-- vault/              # Fichiers .vault chiffres AES-256-GCM
|   `-- security_state.json # Etat des verrous et rate-limits
|-- models/
|   |-- deep_gallery.pkl    # Galerie d'embeddings (pickle)
|   |-- eigenfaces.pkl      # Modele Eigenfaces serialise
|   `-- lbph_model.yml      # Modele LBPH (format OpenCV)
`-- evaluation/
    |-- roc_curve.png
    |-- far_frr_curve.png
    `-- metrics.csv
\end{lstlisting}

% =============================================================================
\section{Revue de Littérature et Concepts Théoriques}
% =============================================================================

\subsection{Détection de visage}

\subsubsection{Cascades de Haar (Viola \& Jones, 2001)}

La méthode de Viola et Jones~\cite{viola2001} est l'un des algorithmes de
détection d'objets les plus influents. Elle repose sur trois innovations :

\begin{enumerate}
  \item \textbf{Image intégrale} : permet de calculer des caractéristiques
        rectangulaires (features de Haar) en temps constant, indépendamment de
        la taille de la fenêtre.
  \item \textbf{AdaBoost} : sélectionne un sous-ensemble de caractéristiques
        discriminantes parmi des milliers de candidats.
  \item \textbf{Cascade de classificateurs} : rejette rapidement les zones sans
        visage dès les premières étapes, réduisant drastiquement le temps de
        calcul.
\end{enumerate}

Dans notre implémentation, le fichier
\texttt{haarcascade\_frontalface\_default.xml} d'OpenCV est chargé via
\texttt{CascadeClassifier.detectMultiScale}, avec un facteur d'échelle de 1.1
et un minimum de 5 voisins.

\subsubsection{Détecteur HOG + SVM (Dlib)}

Le détecteur HOG (\textit{Histogram of Oriented Gradients})~\cite{dalal2005}
calcule, pour chaque cellule d'une image, l'histogramme des orientations des
gradients locaux. Ces histogrammes sont concaténés en un vecteur de
caractéristiques et classifiés par un SVM linéaire entraîné sur des exemples
positifs et négatifs de visages.

Dlib implémente une version optimisée (FHOG) avec un détecteur à fenêtre
glissante multi-échelle. Notre module \texttt{preprocessing.py} utilise
\texttt{dlib.get\_frontal\_face\_detector()}.

\subsubsection{Fusion par NMS}

Pour éviter les détections redondantes, notre pipeline fusionne les résultats
Haar et HOG par \textit{Non-Maximum Suppression} (NMS) avec un seuil IoU de 0,5 :
seule la boîte englobante de plus grande aire est conservée parmi les boîtes
se chevauchant fortement.

\subsection{Points repères faciaux et alignement affine}

\subsubsection{Modèle à 68 points (Dlib)}

Le prédicteur de forme de Dlib~\cite{kazemi2014} utilise une forêt de
régressions d'ensemble (\textit{Ensemble of Regression Trees}) pour localiser
68 points repères anatomiques sur un visage : contour du visage, sourcils, yeux,
nez, bouche.

\begin{infobox}{Indices des points oculaires (68 points)}
  Points 36--41 : œil gauche \quad|\quad Points 42--47 : œil droit
\end{infobox}

\subsubsection{Alignement affine}

À partir des centres géométriques des deux yeux, une transformation affine 2D
(rotation, mise à l'échelle, translation) est calculée pour :
\begin{itemize}
  \item Placer l'œil gauche à 35\,\% et l'œil droit à 65\,\% de la largeur
        de l'image de sortie.
  \item Aligner horizontalement l'axe inter-oculaire.
\end{itemize}

\begin{formulebox}{Transformation affine d'alignement}
\[
  \theta = \arctan\!\left(\frac{y_R - y_L}{x_R - x_L}\right), \quad
  s = \frac{d_{\text{cible}}}{d_{\text{mesure}}}, \quad
  M = \text{getRotationMatrix2D}(\text{centre}, \theta, s)
\]
où $d_{\text{cible}}$ et $d_{\text{mesure}}$ sont respectivement la distance
inter-oculaire désirée et mesurée.
\end{formulebox}

Cette normalisation géométrique améliore considérablement les performances de
tous les modèles en aval en éliminant les variations de pose frontale.

\subsection{Prétraitement et normalisation}

Après alignement, chaque patch facial passe par les étapes suivantes :

\begin{enumerate}
  \item \textbf{Redimensionnement} à $128 \times 128$ pixels (\texttt{INTER\_AREA}
        pour le sous-échantillonnage).
  \item \textbf{Normalisation de l'éclairage} (module \texttt{lighting.py}) :
        compensation des variations d'illumination.
  \item \textbf{Égalisation d'histogramme} (\texttt{cv2.equalizeHist}) :
        redistribution des niveaux de gris pour maximiser le contraste local.
  \item Production de trois variantes : image grise, image couleur, image
        grise normalisée $[0,1]$.
\end{enumerate}

\subsection{Modèles Classiques de Reconnaissance Faciale}

\subsubsection{Eigenfaces – Analyse en Composantes Principales (Turk \& Pentland, 1991)}

Les Eigenfaces~\cite{turk1991} constituent la première approche de reconnaissance
faciale basée sur l'apprentissage statistique. Chaque image faciale
$128 \times 128$ est aplatie en un vecteur de dimension $d = 16\,384$.

\medskip
\textbf{Entraînement :}
\begin{enumerate}
  \item Calcul de la matrice de covariance $C = \frac{1}{n} X^T X$ où $X$
        est la matrice des images centrées.
  \item Décomposition en valeurs propres (PCA via scikit-learn) : conservation
        des $k=150$ premières composantes principales.
  \item Chaque eigenvector (« eigenface ») capture une direction maximale de
        variance dans l'espace des visages.
\end{enumerate}

\textbf{Classification :} projection de la requête dans le sous-espace PCA puis
classification 1-NN (distance euclidienne). Le score de confiance est dérivé de
la distance inverse :
\[
  \text{score} = \frac{1}{1 + d/10}
\]

Notre implémentation utilise \texttt{sklearn.decomposition.PCA(whiten=True)}
dans un \texttt{Pipeline} scikit-learn avec \texttt{StandardScaler} et
\texttt{KNeighborsClassifier}.

\subsubsection{LBPH – Local Binary Pattern Histograms (Ahonen et al., 2006)}

Les LBPH~\cite{ahonen2006} encodent la texture locale d'une image en comparant
chaque pixel à ses voisins selon un motif circulaire de rayon $R$ et $P$
voisins :

\begin{formulebox}{Opérateur LBP}
\[
  \text{LBP}_{P,R}(x_c, y_c) = \sum_{p=0}^{P-1} s(g_p - g_c) \cdot 2^p,
  \quad s(u) = \begin{cases} 1 & u \geq 0 \\ 0 & u < 0 \end{cases}
\]
où $g_c$ est le niveau de gris du pixel central et $g_p$ celui du $p$-ième voisin.
\end{formulebox}

L'image est divisée en une grille de $8 \times 8$ cellules. Pour chaque cellule,
un histogramme de 256 bins est calculé sur les codes LBP, puis concaténé en un
vecteur de description global. La classification se fait par distance $\chi^2$
(intégrée dans le \texttt{LBPHFaceRecognizer} d'OpenCV).

\textbf{Paramètres utilisés :} $R=1$, $P=8$, grille $8\times 8$, seuil de
confiance maximal normalisé à 100.

\subsection{Modèles Profonds de Reconnaissance Faciale}

\subsubsection{FaceNet (Schroff et al., 2015)}

FaceNet~\cite{schroff2015} est un réseau de neurones profond entraîné avec une
\textit{triplet loss} pour produire des embeddings de dimension 128 (ou 512)
dans un espace euclidien tel que :
\begin{itemize}
  \item Les embeddings du \textit{même} individu sont proches ($\|f(x_a) -
        f(x_p)\|_2^2 < \|f(x_a) - f(x_n)\|_2^2 + \alpha$).
  \item Les embeddings d'\textit{individus différents} sont éloignés.
\end{itemize}

\begin{formulebox}{Triplet Loss}
\[
  \mathcal{L} = \sum_i \left[\, \|f(x_i^a) - f(x_i^p)\|_2^2 -
                \|f(x_i^a) - f(x_i^n)\|_2^2 + \alpha \,\right]_+
\]
\end{formulebox}

\subsubsection{ArcFace (Deng et al., 2019)}

ArcFace~\cite{deng2019} améliore FaceNet en ajoutant une marge angulaire $m$
directement dans l'espace des angles entre embeddings et poids du classifieur :

\begin{formulebox}{ArcFace Loss}
\[
  \mathcal{L} = -\frac{1}{N}\sum_{i=1}^{N} \log
  \frac{e^{s(\cos(\theta_{y_i} + m))}}
       {e^{s(\cos(\theta_{y_i} + m))} + \sum_{j \neq y_i} e^{s \cos\theta_j}}
\]
\end{formulebox}

Cette formulation produit des embeddings avec une meilleure séparabilité
inter-classe et une plus grande compacité intra-classe que FaceNet.

\subsubsection{Extraction et comparaison d'embeddings}

Dans notre système, la bibliothèque \textbf{DeepFace} est utilisée comme
abstraction sur TensorFlow :

\begin{lstlisting}[caption={Extraction d'embedding (deep\_model.py)}]
result = DeepFace.represent(
    img_path=face_img,
    model_name=self.backend,   # "Facenet" ou "ArcFace"
    enforce_detection=False,
    detector_backend="skip",   # pre-cropped par notre pipeline
)
embedding = np.array(result[0]["embedding"], dtype=np.float32)
embedding /= np.linalg.norm(embedding)  # normalisation L2
\end{lstlisting}

La \textbf{similarité cosinus} est utilisée pour la comparaison :
\[
  \text{sim}(\mathbf{a}, \mathbf{b}) = 1 - \frac{\mathbf{a} \cdot \mathbf{b}}
  {\|\mathbf{a}\|\,\|\mathbf{b}\|}
  \quad\in [0, 1]
\]

Un seuil de distance cosinus de 0.40 est appliqué : si la distance est
inférieure au seuil, les deux visages sont considérés comme appartenant à la
même identité.

\subsection{Fusion de Scores}

\subsubsection{Motivation}

Chaque modèle capture des caractéristiques différentes :
\begin{itemize}
  \item Les modèles classiques sont robustes, légers et interprétables.
  \item Les modèles profonds sont plus discriminants et généralisent mieux.
\end{itemize}

La \textbf{fusion de scores} combine les deux en exploitant leur
complémentarité.

\subsubsection{Fusion linéaire pondérée}

\begin{formulebox}{Score fusionné}
\[
  s_{\text{fused}} = w_c \cdot s_{\text{classical}} + w_d \cdot s_{\text{deep}},
  \quad w_c + w_d = 1
\]
\[
  w_c = 0.4, \quad w_d = 0.6, \quad \text{seuil de décision} = 0.6
\]
\end{formulebox}

Une pénalité de 50\,\% est appliquée sur le score fusionné si l'identité
revendiquée ne correspond pas à celle retournée par le modèle profond.

\subsection{Détection de Vivacité (\textit{Liveness Detection})}

La détection de vivacité vise à distinguer un visage réel d'une attaque par
présentation (\textit{Presentation Attack}) : photo imprimée, vidéo rejouée,
masque 3D.

\subsubsection{Eye Aspect Ratio (EAR) – Détection de clignements}

L'EAR~\cite{soukupova2016} mesure le rapport entre la hauteur et la largeur
d'un œil à partir des 6 points repères oculaires :

\begin{formulebox}{Eye Aspect Ratio}
\[
  \text{EAR} = \frac{\|p_2 - p_6\| + \|p_3 - p_5\|}{2 \cdot \|p_1 - p_4\|}
\]
\end{formulebox}

Lors d'un clignement, l'EAR chute brusquement puis remonte. Une photo ou un
écran ne peut pas reproduire ce comportement temporel. Notre détecteur exige
\textbf{1 clignement confirmé} sur au moins 3 frames consécutives avec EAR $<$
0.25 pour valider la vivacité.

\subsubsection{Analyse de texture LBP}

La texture d'un vrai visage est plus riche et variée que celle d'une impression
sur papier ou d'un écran. Un descripteur LBP en grille $3\times 3$ est calculé
(vecteur de $9 \times 256 = 2304$ dimensions) et analysé :
\begin{itemize}
  \item Si un classifieur sklearn entraîné est disponible, \texttt{predict\_proba}
        donne une probabilité de vivacité.
  \item Sinon, une heuristique basée sur la \textbf{variance de l'histogramme}
        LBP est utilisée : plus la variance est élevée, plus le visage est
        probablement réel.
\end{itemize}

\subsection{Détection de Deepfakes par Analyse Spectrale}

\subsubsection{Artefacts des réseaux GAN}

Les visages générés par GAN (\textit{Generative Adversarial Network}) présentent
des artefacts hautes fréquences liés au sur-échantillonnage par convolution
transposée : des motifs en damier (\textit{checkerboard artifacts}) se
manifestent dans le spectre fréquentiel~\cite{li2021}.

\subsubsection{Pipeline FFT}

Notre détecteur (\texttt{deepfake\_detector.py}) calcule trois métriques sur
le spectre log-Fourier 2D de l'image :

\begin{longtable}{>{\bfseries}p{4.5cm} p{9cm}}
\toprule
Métrique & Définition \\
\midrule
\endhead
\bottomrule
\endfoot
Ratio hautes fréquences & Fraction de l'énergie spectrale en dehors d'un
                          disque central de rayon 30\,\% (GAN $\Rightarrow$
                          plus d'HF). \\[4pt]
Score checkerboard & Énergie moyenne aux coins de Nyquist de l'image
                     (artefacts de sur-échantillonnage GAN). \\[4pt]
Platitude spectrale & Inverse de la variance radiale du spectre
                      (vraies images suivent une loi en $1/f$,
                       les GAN produisent un spectre plus plat). \\
\end{longtable}

\begin{formulebox}{Score de fausseté}
\[
  s_{\text{fake}} = 0.4 \cdot s_{\text{HF}} + 0.4 \cdot s_{\text{CB}} +
                    0.2 \cdot s_{\text{flat}}
\]
Si $s_{\text{fake}} \geq 0.55$, le visage est classé comme synthétique.
\end{formulebox}

\subsection{Biométrie Annulable (\textit{Cancelable Biometrics})}

\subsubsection{Problématique}

Les templates biométriques sont permanents : si un embedding facial est
compromis, l'individu ne peut pas ``renouveler'' son visage. La
\textbf{biométrie annulable}~\cite{ratha2001} répond à ce problème en
transformant le template de manière non-inversible mais reproductible.

\subsubsection{BioHashing par projection orthogonale}

Notre implémentation suit le paradigme BioHashing~\cite{jin2004} :

\begin{enumerate}
  \item Un \textbf{token secret} (PIN, carte à puce) est haché par SHA-256
        pour obtenir une graine déterministe.
  \item Une matrice de projection orthogonale aléatoire $P \in \mathbb{R}^{256
        \times 512}$ est générée par décomposition QR d'une matrice gaussienne
        secrète.
  \item L'embedding $\mathbf{e} \in \mathbb{R}^{512}$ est projeté :
        $\mathbf{v} = P\mathbf{e} \in \mathbb{R}^{256}$.
  \item Binarisation par signe : $b_i = \mathbf{1}[v_i \geq 0]$.
\end{enumerate}

\begin{formulebox}{Template protégé}
\[
  \mathbf{t} = \text{sign}(P_{\text{token}} \cdot \mathbf{e}) \in \{0,1\}^{256}
\]
\end{formulebox}

\textbf{Propriétés garanties :}
\begin{itemize}
  \item \textbf{Non-inversibilité} : sans le token, $\mathbf{e}$ ne peut pas
        être récupéré depuis $\mathbf{t}$.
  \item \textbf{Révocabilité} : un nouveau token génère un template
        décorrélé du précédent.
  \item \textbf{Diversité} : différents tokens $\Rightarrow$ différents templates
        $\Rightarrow$ impossible de relier deux systèmes.
\end{itemize}

La comparaison entre deux templates binaires utilise la
\textbf{similarité de Hamming normalisée} :
\[
  \text{sim}_H(\mathbf{t}_1, \mathbf{t}_2) = \frac{1}{n} \sum_{i=1}^n
  \mathbf{1}[t_{1,i} = t_{2,i}]
\]

\subsection{Coffre-Fort Chiffré (AES-256-GCM)}

\subsubsection{Chiffrement authentifié}

Les templates protégés sont stockés chiffrés avec
\textbf{AES-256-GCM}~\cite{nist2007} (AEAD : Authenticated Encryption with
Associated Data), qui garantit à la fois la confidentialité et l'intégrité.

\subsubsection{Dérivation de clé PBKDF2}

La clé AES de 256 bits est dérivée du mot de passe utilisateur via
\textbf{PBKDF2-HMAC-SHA256}~\cite{rfc8018} avec 100\,000 itérations et un
sel aléatoire de 16 octets :

\begin{formulebox}{Dérivation de clé}
\[
  k = \text{PBKDF2}(\text{password}, \text{salt}, \text{iter}=100\,000,
  \text{len}=32, \text{PRF}=\text{HMAC-SHA256})
\]
\end{formulebox}

\subsubsection{Format de fichier}

Chaque utilisateur dispose d'un fichier \texttt{.vault} dont le format binaire
est :
\[
  \underbrace{\text{salt}}_{\text{16 B}} \;\|\;
  \underbrace{\text{nonce}}_{\text{12 B}} \;\|\;
  \underbrace{\text{ciphertext + GCM tag}}_{\text{variable}}
\]

\subsection{Moniteur de Sécurité}

Le module \texttt{security\_monitor.py} implémente deux politiques de sécurité
complémentaires :

\begin{enumerate}
  \item \textbf{Limitation de débit (sliding window)} : au plus 3 tentatives
        par fenêtre glissante de 30 secondes. La liste des timestamps est
        filtrée à chaque contrôle.
  \item \textbf{Verrouillage progressif} : après 5 tentatives consécutives
        échouées, le compte est verrouillé pendant 5 minutes. L'état survit
        aux redémarrages grâce à une persistance JSON.
\end{enumerate}

\subsection{Évaluation des Systèmes Biométriques}

Les métriques standard ISO/IEC 19795 pour l'évaluation biométrique sont
implémentées dans \texttt{evaluation.py} :

\begin{longtable}{>{\bfseries}p{4cm} p{9.5cm}}
\toprule
Métrique & Définition \\
\midrule
\endhead
\bottomrule
\endfoot
FAR (\textit{False Accept Rate}) & Taux de faux positifs : fraction des
  comparaisons imposteur acceptées. $\text{FAR} = \text{FP}/(\text{FP}+\text{TN})$ \\[4pt]
FRR (\textit{False Reject Rate}) & Taux de faux négatifs : fraction des
  comparaisons authentiques rejetées. $\text{FRR} = \text{FN}/(\text{FN}+\text{TP})$ \\[4pt]
EER (\textit{Equal Error Rate}) & Seuil où FAR $=$ FRR ; indicateur scalaire
  de performance globale. \\[4pt]
Courbe ROC & TPR en fonction de FPR pour tous les seuils ; aire sous la
  courbe (AUC) proche de 1 $\Rightarrow$ meilleur système. \\
\end{longtable}

% =============================================================================
\section{Flux d'Authentification}
% =============================================================================

\subsection{Inscription (\textit{Enrollment})}

\begin{enumerate}
  \item L'utilisateur saisit son prénom, nom et date de naissance.
  \item 10 images sont capturées via la webcam (ou fournies comme fichiers).
  \item Pour chaque image : détection, alignement, normalisation.
  \item Extraction de l'embedding FaceNet/ArcFace moyen sur les 10 images.
  \item L'embedding est projeté via BioHashing (token secret de l'utilisateur).
  \item Le template protégé est chiffré AES-256-GCM et stocké dans le vault.
  \item L'embedding brut est stocké dans la galerie pickle pour l'identification.
  \item Les recadrages faciaux sont sauvegardés pour le réentraînement classique.
  \item Le profil (nom, dob) est enregistré dans \texttt{users.json}.
\end{enumerate}

\subsection{Authentification}

\begin{enumerate}
  \item Le moniteur de sécurité vérifie l'absence de verrouillage et de
        rate-limiting.
  \item Vivacité : clignement requis (EAR) sur la flux vidéo en temps réel.
  \item Détection de deepfake (FFT) sur chaque frame.
  \item Preprocessing : détection, alignement, normalisation.
  \item \textbf{Pipeline profond} : identification dans la galerie par
        similarité cosinus $\Rightarrow$ identité + score $s_d$.
  \item \textbf{Pipeline classique} (si modèles entraînés) : LBPH ou
        Eigenfaces $\Rightarrow$ score $s_c$.
  \item Fusion : $s_f = 0.4 \, s_c + 0.6 \, s_d$.
  \item Décision : GRANTED si $s_f \geq 0.6$, DENIED sinon.
  \item Journalisation et mise à jour du moniteur de sécurité.
\end{enumerate}

Le seuil de confirmation (\textit{streak}) est fixé à 3 frames consécutives
positives pour réduire les faux positifs transitoires.

% =============================================================================
\section{Interface Utilisateur}
% =============================================================================

L'interface est construite avec \textbf{Tkinter}, la bibliothèque graphique
standard de Python, sans dépendance externe supplémentaire.

\subsection{Écrans}

\begin{longtable}{>{\bfseries}p{3.5cm} p{10cm}}
\toprule
Écran & Description \\
\midrule
\endhead
\bottomrule
\endfoot
HomeScreen & Écran d'accueil avec boutons ``S'identifier'' et
             ``Créer un compte''. Palette sombre, typographie minimaliste. \\[4pt]
RegisterScreen & Formulaire de saisie du profil utilisateur + capture
                 webcam guidée. Indicateur de progression des captures. \\[4pt]
LoginScreen & Flux vidéo temps réel avec overlay de statut, barre de
              confirmation progressive et indicateurs de vivacité. \\[4pt]
WelcomeScreen & Écran de succès post-authentification : coche verte, nom
                complet, date de naissance, bouton de déconnexion. \\
\end{longtable}

\subsection{Design}

Le design applique un thème sombre cohérent (\texttt{BG = \#0d0d0d},
\texttt{TEXT = \#f0f0f0}) avec une typographie épurée. Les composants
réutilisables (\texttt{\_lbl}, \texttt{\_btn}) centralisent la cohérence visuelle.
Un système de confirmation par « streak » (3 frames consécutives positives)
avec barre de progression visuelle améliore l'expérience utilisateur.

% =============================================================================
\section{Tests}
% =============================================================================

Le projet dispose d'une suite de tests unitaires et d'intégration (\texttt{pytest}) :

\begin{longtable}{>{\bfseries}p{5cm} p{9cm}}
\toprule
Fichier de test & Couverture \\
\midrule
\endhead
\bottomrule
\endfoot
test\_authenticator.py & Fusion de scores, pénalité d'identité, seuils. \\[3pt]
test\_cancelable\_biometrics.py & Reproductibilité de la projection, révocabilité,
                                   similarité de Hamming. \\[3pt]
test\_evaluation.py & Calcul FAR, FRR, EER sur des jeux de données synthétiques. \\[3pt]
test\_security\_monitor.py & Verrouillage progressif, rate-limiting, persistance. \\[3pt]
test\_deepfake\_detector.py & Score FFT sur images synthétiques vs réelles. \\[3pt]
test\_liveness.py & EAR blink detection, heuristique LBP texture. \\[3pt]
test\_vault.py & Chiffrement/déchiffrement AES-GCM, résistance aux mauvais
                 mots de passe. \\
\end{longtable}

% =============================================================================
\section{Sécurité et Vie Privée}
% =============================================================================

\subsection{Modèle de menace}

Le système est conçu pour résister aux attaques suivantes :

\begin{itemize}
  \item \textbf{Présentation d'artefact} (photo, vidéo) : détecté par EAR.
  \item \textbf{Deepfake} : détecté par analyse spectrale FFT.
  \item \textbf{Usurpation de texture LBP} : contrecarré par la vérification
        de texture.
  \item \textbf{Vol de base de données} : les templates sont chiffrés
        AES-256-GCM ; les embeddings bruts ne permettent pas de reconstruire
        les images originales.
  \item \textbf{Réutilisation de template compromis} : la biométrie annulable
        permet l'émission d'un nouveau token.
  \item \textbf{Brute force} : le moniteur de sécurité limite les tentatives
        et impose des verrouillages.
  \item \textbf{Liaison inter-systèmes} : des tokens différents génèrent des
        templates décorrélés sur différents systèmes.
\end{itemize}

\subsection{Conformité RGPD}

La présence d'un fichier de consentement (\texttt{consent.json}) et l'absence
de stockage d'images biométriques brutes en clair (seuls les embeddings
chiffrés sont conservés) s'inscrivent dans les principes de minimisation et de
protection by design du RGPD.

% =============================================================================
\section{Conclusion}
% =============================================================================

Ce projet démontre la faisabilité d'un système d'authentification faciale
complet, sécurisé et respectueux de la vie privée, intégrant :

\begin{itemize}
  \item Une \textbf{double approche} (classique + profonde) avec fusion de
        scores pour une robustesse accrue.
  \item Une \textbf{chaîne anti-spoof} multicouche (vivacité temporelle, texture
        spatiale, analyse spectrale deepfake).
  \item Une \textbf{protection cryptographique} des données biométriques
        (AES-256-GCM, PBKDF2, BioHashing révocable).
  \item Un \textbf{moniteur de sécurité} résistant aux attaques par force brute.
  \item Une \textbf{interface graphique} temps réel intégrée.
\end{itemize}

Les perspectives d'amélioration incluent l'intégration d'une détection de
vivacité 3D (estimation de la profondeur), l'utilisation de modèles plus
récents (ArcFace R100, ElasticFace), et le déploiement sur une architecture
client-serveur avec authentification multi-facteurs.

% =============================================================================
% Bibliographie
% =============================================================================
\newpage
\begin{thebibliography}{99}

\bibitem{viola2001}
P. Viola and M. Jones, ``Rapid object detection using a boosted cascade of
simple features,'' \textit{Proc. CVPR}, 2001.

\bibitem{dalal2005}
N. Dalal and B. Triggs, ``Histograms of oriented gradients for human
detection,'' \textit{Proc. CVPR}, 2005.

\bibitem{kazemi2014}
V. Kazemi and J. Sullivan, ``One millisecond face alignment with an ensemble
of regression trees,'' \textit{Proc. CVPR}, 2014.

\bibitem{turk1991}
M. Turk and A. Pentland, ``Eigenfaces for recognition,''
\textit{Journal of Cognitive Neuroscience}, vol.~3, no.~1, 1991.

\bibitem{ahonen2006}
T. Ahonen, A. Hadid, and M. Pietikäinen, ``Face description with local binary
patterns: Application to face recognition,'' \textit{IEEE TPAMI},
vol.~28, no.~12, 2006.

\bibitem{schroff2015}
F. Schroff, D. Kalenichenko, and J. Philbin, ``FaceNet: A unified embedding
for face recognition and clustering,'' \textit{Proc. CVPR}, 2015.

\bibitem{deng2019}
J. Deng, J. Guo, N. Xue, and S. Zafeiriou, ``ArcFace: Additive angular
margin loss for deep face recognition,'' \textit{Proc. CVPR}, 2019.

\bibitem{soukupova2016}
T. Soukupova and J. Cech, ``Real-time eye blink detection using facial
landmarks,'' \textit{Computer Vision Winter Workshop}, 2016.

\bibitem{li2021}
L. Li \textit{et al.}, ``Thinking in frequency: Face forgery detection by
mining frequency-aware clues,'' \textit{Proc. ECCV}, 2020.

\bibitem{ratha2001}
N. K. Ratha, J. H. Connell, and R. M. Bolle, ``Enhancing security and privacy
in biometrics-based authentication systems,'' \textit{IBM Systems Journal},
vol.~40, no.~3, 2001.

\bibitem{jin2004}
A. T. B. Jin, D. N. C. Ling, and A. Goh, ``Biohashing: Two factor
authentication featuring fingerprint data and tokenised random number,''
\textit{Pattern Recognition}, vol.~37, no.~11, 2004.

\bibitem{nist2007}
M. Dworkin, ``Recommendation for Block Cipher Modes of Operation:
Galois/Counter Mode (GCM) and GMAC,'' NIST SP 800-38D, 2007.

\bibitem{rfc8018}
K. Moriarty, B. Kaliski, and A. Rusch, ``PKCS \#5: Password-Based
Cryptography Specification Version 2.1,'' RFC 8018, 2017.

\end{thebibliography}

\end{document}
